%% LyX 2.4.4 created this file.  For more info, see https://www.lyx.org/.
%% Do not edit unless you really know what you are doing.
\documentclass[english]{article}
\usepackage[T1]{fontenc}
\usepackage[utf8]{inputenc}
\setcounter{secnumdepth}{-2}
\setlength{\parindent}{0bp}
\usepackage{amsmath}
\usepackage{amsthm}

\makeatletter
%%%%%%%%%%%%%%%%%%%%%%%%%%%%%% Textclass specific LaTeX commands.
\theoremstyle{definition}
\newtheorem*{problem*}{\protect\problemname}

\@ifundefined{date}{}{\date{}}
%%%%%%%%%%%%%%%%%%%%%%%%%%%%%% User specified LaTeX commands.
\usepackage{graphicx}
\usepackage{geometry}
\newcommand{\limi}{\underset{n\rightarrow\infty}{\lim}}
\newcommand{\limxz}{\underset{x\rightarrow x_{0}}{\lim}}
\newcommand{\limfxz}{\underset{x\rightarrow x_{0}}{\lim}f(x)}
\newcommand{\limfxm}{\underset{x\rightarrow x_{0}^{-}}{\lim}f(x)}
\newcommand{\limfxp}{\underset{x\rightarrow x_{0}^{+}}{\lim}f(x)}
\newcommand{\limxm}{\underset{x\rightarrow x_{0}^{-}}{\lim}}
\newcommand{\limxp}{\underset{x\rightarrow x_{0}^{+}}{\lim}}
\newcommand{\sqrm}{M_{m\times m}(\mathbb{F})}
\newcommand{\seqa}{(a_{n})_{n=1}^{\infty}}
\newcommand{\seqx}{(x_{n})_{n=1}^{\infty}}
\newcommand{\funcdr}{f:D\rightarrow\mathbb{R}}
\newcommand{\der}{\limxz\frac{f(x)-f(x_{0})}{x-x_{0}}}
\usepackage[all]{pdfprivacy}

\makeatother

\usepackage{babel}
\providecommand{\problemname}{Problem}
\begin{document}
\title{{\Large Semester 2026A, Exam A, Q4\vspace{-2em}}}

\maketitle
\begin{problem*}
Let $V$ be a vector space and let $v_{1},\dots,v_{n},u,w\in V$
be such that the sequence $(v_{1},v_{2},\dots,v_{n},u)$ is linearly
dependent, while the sequence $(v_{1},\dots,v_{n},w)$ is linearly
independent. 

Is it necessarily the case that $w\notin\operatorname{Span}(v_{1},v_{2},\dots,v_{n},u)$?
\end{problem*}

\medskip{}

\rule[0.5ex]{1\columnwidth}{1pt}

\medskip{}

\begin{proof}
The sequence $(v_{1},\dots,v_{n},w)$ is linearly independent, which
means that the sequence $(v_{1},\dots,v_{n})$ is also linearly independent
(recall that if a sequence is linearly dependent one of the vectors
is a linear combination of the others, hence a sequence that is already
linearly dependent will stay linearly dependent when we add a vector
to it). 

Recall that if you add a vector to a linearly independent sequence
which does not belong to the span of this sequence, you receive a
new sequence which is still linearly independent, hence since $(v_{1},v_{2},\dots,v_{n},u)$
is linearly dependent we get that 
\[
u\in\operatorname{Span}(v_{1},v_{2},\dots,v_{n})
\]

Recall that if a vector belongs to the span of a sequence, the span
of this sequence is equal to the span of this sequence along with
said vector, hence 

\[
\operatorname{Span}(v_{1},v_{2},\dots,v_{n})=\operatorname{Span}(v_{1},v_{2},\dots,v_{n},u)
\]
Which means $w$ cannot belong to $\operatorname{Span}(v_{1},v_{2},\dots,v_{n},u)$,
because in that 
\[
w\in\operatorname{Span}(v_{1},v_{2},\dots,v_{n})
\]

which would mean the sequence $(v_{1},\dots,v_{n},w)$ is linearly
dependent.
\end{proof}
\begin{quote}
This is an exam problem I translated and solved. The original exam
was written by the Linear Algebra 1 course lecturers at HUJI.
\end{quote}

\end{document}
